\documentclass[11pt, a4paper, twoside, titlepage]{scrbook}
\usepackage{ccicons}
\usepackage{fontawesome5} % Required for the cover watermark

% ==========================================
% KOMA-SCRIPT FONT SETTINGS & PREMIUM FONTS
% ==========================================
\usepackage{newpxtext, newpxmath}       % Palatino font for body and math
\usepackage[scaled=0.95]{inconsolata}   % Sleek modern monospace font for code
\usepackage[scaled=0.92]{helvet}       % Clean sans-serif for headings

\setkomafont{disposition}{\sffamily\bfseries}
\setkomafont{descriptionlabel}{\normalfont\bfseries}
\setkomafont{pageheadfoot}{\normalfont\small}

% ==========================================
% PACKAGES
% ==========================================
\usepackage[utf8]{inputenc}
\usepackage[T1]{fontenc}
\usepackage[english]{babel}
\usepackage{etoolbox} % Required for advanced hooks and counters

% --- PREVENT \Bbbk AND \openbox COLLISIONS ---
\let\Bbbk\undefined
\let\openbox\undefined
% ---------------------------------------------

% Matematica e Informatica Teorica
\usepackage{amsmath, amssymb, amsthm, mathtools, cancel}
\usepackage{stmaryrd} % Simboli aggiuntivi (es. parentesi per semantica)

% Layout e Geometria
\usepackage{geometry}
\geometry{a4paper, inner=2.5cm, outer=2.5cm, bindingoffset=0.5cm, top=2.5cm, bottom=2.5cm}
\usepackage{multicol} % For 2-column acronyms

% Grafica e Colori
\usepackage{graphicx}
\usepackage{xcolor}
\usepackage{tikz} % Per grafi, automi, architetture
\usetikzlibrary{automata, positioning, arrows, arrows.meta}

% Box per Teoremi, Definizioni, Algoritmi
\usepackage{tcolorbox}
\tcbuselibrary{theorems, breakable, skins}

% Codice e Pseudocodice
\usepackage{listings}
\usepackage{algorithm}
\usepackage{algpseudocode}

% ==========================================
% CUSTOMIZATION & BRANDING (UNITN STYLE)
% ==========================================
\definecolor{unitnred}{RGB}{178, 35, 52}
\definecolor{codebg}{RGB}{245, 245, 245}

% ==========================================
% MODERN CHAPTER HEADINGS
% ==========================================
\makeatletter
\renewcommand*{\chapterlinesformat}[3]{%
  \Ifstr{#1}{chapter}{%
    \parbox[b]{\textwidth}{%
      \raggedright
      \Ifstr{#2}{}{}{% If there is a chapter number
        {\fontsize{70}{75}\selectfont\color{unitnred!20}#2}\\[-2ex]
      }%
      \Huge\bfseries\sffamily\color{black}#3\par\vspace{1cm}
    }%
  }{% For unnumbered chapters (like Preface/Appendix)
    \@hangfrom{#2}{#3}%
  }%
}
\makeatother

% Link e Riferimenti
\usepackage{hyperref}
\hypersetup{
    colorlinks=true,
    linkcolor=unitnred,
    filecolor=magenta,      
    urlcolor=unitnred,
    citecolor=blue,
    pdftitle={Lecture Notes - Computability and Computational Complexity},
    pdfauthor={Luca Facchini},
    pdfsubject={Master's Degree in Computer Science - University of Trento},
}

\usepackage[acronym,toc]{glossaries-extra}
\makeglossaries

% --- AUTOMATED ACRONYM USAGE COUNTER HOOKS ---
\makeatletter
\appto\glslinkpostsetkeys{%
    \ifcsundef{acrocount@\glslabel}{%
        \csgdef{acrocount@\glslabel}{1}%
    }{%
        \csxdef{acrocount@\glslabel}{\the\numexpr\csuse{acrocount@\glslabel} + 1\relax}%
    }%
}

\renewcommand{\glsxtrpostdescacronym}{%
    \ifcsdef{acrocount@\glscurrententrylabel}{%
        \ifnum\csuse{acrocount@\glscurrententrylabel}=1\relax
            \space \textit{[1]}%
        \else
            \space \textit{[\csuse{acrocount@\glscurrententrylabel}]}%
        \fi
    }{}%
}
\makeatother

% Clever references
\usepackage{cleveref}
\newcommand{\vertmag}{\rotatebox{90}{$\,<$}}
\newcommand{\equalto}[2]{\underset{\scriptstyle\overset{\mkern4mu\vertmag}{#2}}{#1}}

% ==========================================
% ACRONYMS DEFINITIONS
% ==========================================
\newacronym{cs}{CS}{Computer Science}
\newacronym{pm}{PM}{\textit{Principia Mathematica}}
\newacronym{zf}{ZF}{Zermelo-Fraenkel set theory}

% Stile del Codice
\lstset{
    backgroundcolor=\color{codebg},
    basicstyle=\ttfamily\footnotesize,
    keywordstyle=\color{unitnred}\bfseries,
    commentstyle=\color{green!60!black}\itshape,
    stringstyle=\color{blue},
    breaklines=true,
    frame=single,
    rulecolor=\color{lightgray},
    numbers=left,
    numberstyle=\tiny\color{gray},
    showstringspaces=false
}

% ==========================================
% TCOLORBOX THEOREMS & DEFINITIONS
% ==========================================
\newtcbtheorem[number within=chapter,
  crefname={definition}{definitions},
  Crefname={Definition}{Definitions},
  list inside=deflist
]{definition}{Definition}{%
  colback=unitnred!5,colframe=unitnred!80!black,fonttitle=\bfseries, breakable%
}{def}

\newtcbtheorem[number within=chapter,
  crefname={theorem}{theorems},
  Crefname={Theorem}{Theorems},
  list inside=thmlist
]{theorem}{Theorem}{%
  colback=blue!5,colframe=blue!80!black,fonttitle=\bfseries, breakable%
}{thm}

\newtcbtheorem[number within=chapter,
  crefname={lemma}{lemmas},
  Crefname={Lemma}{Lemmas},
  list inside=lemlist
]{lemma}{Lemma}{%
  colback=cyan!5,colframe=cyan!80!black,fonttitle=\bfseries, breakable%
}{lem}

\newtcbtheorem[number within=chapter,
  crefname={observation}{observations},
  Crefname={Observation}{Observations},
  list inside=obslist
]{observation}{Observation}{%
  colback=orange!5,colframe=orange!80!black,fonttitle=\bfseries, breakable%
}{obs}

\newtcbtheorem[number within=chapter,
  crefname={concept}{concepts},
  Crefname={Concept}{Concepts},
  list inside=cptlist
]{concept}{Concept}{%
  colback=green!5!white,colframe=green!60!black,fonttitle=\bfseries, breakable%
}{cpt}

\newtcbtheorem[number within=chapter,
  crefname={axiom}{axioms},
  Crefname={Axiom}{Axioms},
  list inside=axlist
]{axiom}{Axiom}{%
  colback=purple!5,colframe=purple!80!black,fonttitle=\bfseries, breakable%
}{ax}

% ==========================================
% HEADER & FOOTER (KOMA-Script Native)
% ==========================================
\usepackage[headsepline=0.4pt, footsepline=0.4pt, plainfootsepline]{scrlayer-scrpage}
\clearpairofpagestyles

\automark[section]{chapter}
\ohead{\headmark} 
\ihead{\textbf{Computability and Computational Complexity}} 

\ifoot[A.Y. 2026/2027 -- \textit{Notes by Luca Facchini (Matr. 273741)}]{A.Y. 2026/2027 -- \textit{Notes by Luca Facchini (Matr. 273741)}}
\ofoot[Page \pagemark]{Page \pagemark}

\addtokomafont{headsepline}{\color{unitnred}}
\addtokomafont{footsepline}{\color{unitnred}}
\setlength{\headheight}{24pt}

\newcommand{\lecturestop}[1]{%
    \par\vspace{1em}
    \begin{tcolorbox}[
        colback=gray!10,
        colframe=gray!60,
        coltext=gray!60!black,
        fontupper=\centering\bfseries\sffamily,
        boxrule=0.5pt,
        arc=2mm,
        left=4mm, right=4mm, top=2mm, bottom=2mm
    ]
        Lecture stopped here on #1.
    \end{tcolorbox}
    \par\vspace{1em}
}

% ==========================================
% DOCUMENT BEGINS
% ==========================================
\begin{document}
\frontmatter

% --- TITLE PAGE ---
\begin{titlepage}
    % Full-page background graphics using Ti*k*Z
    \begin{tikzpicture}[remember picture, overlay]
        % 1. Bold UniTn red vertical bar on the left edge
        \fill[unitnred] (current page.south west) rectangle ([xshift=1.5cm]current page.north west);
        % 2. Subtle gray geometric accent at the bottom right
        \fill[gray!10] (current page.south east) -- ([yshift=6cm]current page.south east) -- ([xshift=-8cm]current page.south east) -- cycle;
        % 3. Giant faint infinity watermark icon in background
        \node[text=gray!10, font=\fontsize{250}{250}\selectfont, anchor=center] at (current page.center) {\faInfinity};
    \end{tikzpicture}
    
    % Push content down and to the right to clear the red bar
    \vspace*{2cm}
    \noindent\hspace{1cm}%
    \begin{minipage}{\dimexpr\textwidth-1cm\relax}
        
        % University Logo
        \includegraphics[width=4.5cm]{UniTn.jpg} 
        
        \vspace{4cm}
        
        % Main Title Area
        {\Huge \bfseries \sffamily \color{unitnred} Computability and Computational Complexity\par}
        \vspace{0.5cm}
        {\LARGE \itshape Lecture Notes\par}
        \vspace{0.5cm}
        {\Large A.Y. 2026/2027\par}
        
        \vfill
        
        % Author Info Block
        \begin{flushleft}
            {\Large \textbf{Luca Facchini} (Matr. 273741)}\\[0.2cm]
            \textit{Master's Degree in Computer Science}\\
            University of Trento\\[0.3cm]
            \href{mailto:luca.facchini-1@studenti.unitn.it}{\texttt{luca.facchini-1@studenti.unitn.it}}\\
            \href{mailto:luca@fc-software.it}{\texttt{luca@fc-software.it}}
        \end{flushleft}
        
        \vspace{1.5cm}
        
        % Footer line and date
        \textcolor{gray}{\rule{\linewidth}{0.8pt}}
        \vspace{0.2cm}\\
        {\small \textbf{Last updated:} \today \hfill \textit{University of Trento}}
        
    \end{minipage}
\end{titlepage}
\clearpage

% --- COPYRIGHT / LEGAL PAGE (Back of Title Page) ---
\newpage
\thispagestyle{empty}
\vspace*{\fill} % Pushes the text to the bottom of the page

\noindent \textbf{About this Document}\\
These notes cover the topics of the \textbf{Computability and Computational Complexity} course, held by Prof. Mauro Brunato during the Master's Degree in Computer Science at the University of Trento. The content herein is derived from class lectures, course recordings, and the professor's official slides and notes.

\vspace{1em}
\noindent \textbf{License \& Copyright}\\
\textcopyright\ 2026 Luca Facchini. My personal contributions, transcriptions, and formatting within this document are distributed under the Creative Commons Attribution-NonCommercial 4.0 International license (\textbf{CC BY-NC 4.0} \ccbynceu). The full text of the license is available in the \texttt{LICENSE} file at the root of this project. Please note that any original intellectual property, figures, or direct excerpts sourced from the professor remain strictly under their respective original copyrights and licenses.

\vspace{1em}
\noindent \textbf{Disclaimer}\\
This material is provided ``AS IS'', without warranty of any kind, express or implied, including but not limited to the warranties of accuracy, completeness, or fitness for a particular purpose. These are independently produced student notes and do not constitute an official university publication. 

\vspace{1em}
\noindent \textit{If you find any typos or errors, please contact me at the email on the title page.}
\vspace{2cm}
\cleardoublepage

% --- TABLE OF CONTENTS ---
\addtokomafont{chapterentry}{\color{unitnred}}
\addtokomafont{chapterentrypagenumber}{\color{unitnred}}
\tableofcontents
\cleardoublepage
\mainmatter

% --- PREFACE / SCOPE ---
\addchap{Preface: Scope and Purpose}

\noindent \textbf{The Motivation}\\
These notes were initially born out of a personal need: the desire to organize, synthesize, and rigorously typeset the vast amount of information presented during the lectures. Writing and structuring this material in \LaTeX\ served as my primary method of active recall and study. However, recognizing how challenging it can be to study from scattered slides, recordings, and fragmented notes, I decided to open-source this project. The goal is to provide current and future students of the Master's Degree in Computer Science with a unified, structured, and accessible study companion.

\vspace{1em}
\noindent \textbf{The Scope}\\
The scope of this document is to cover the core theoretical concepts, definitions, mathematical proofs, and computational models discussed throughout the course. It is designed to act as a comprehensive bridge between the professor's official slides and the verbal explanations provided during class. 

\vspace{1em}
\noindent \textbf{How to Use This Material}\\
While I have made every effort to ensure the mathematical and technical rigor of this document, it is important to state what these notes are \textit{not}. They are \textbf{not a replacement} for the official recommended textbooks, nor are they a substitute for attending lectures and engaging with the professor. They are a student-made study guide. You should use them to clarify complex topics, review for exams, and quickly reference formal definitions.

\vspace{1em}
\noindent \textbf{Collaborative Spirit}\\
Because this is an actively maintained, open-source project, it is subject to continuous improvement. If you are studying from this document and spot a typographical error, a mathematical inaccuracy, or a missing concept, please do not hesitate to reach out via the contact information on the title page. Collaborative corrections are highly encouraged and deeply appreciated.

\vspace{2em}
\noindent \textit{Happy studying,}\\
\textbf{Luca Facchini}\\
BSc Graduate in Computer Science

\cleardoublepage

% ==========================================
% LECTURES BEGIN
% ==========================================
\chapter{A Short History of Computability [14/09/2026]}

Leibniz (1646-1716) made an attempt to formalize human reasoning, whose motto was ``Calculemus!''. The criteria were:
\begin{description}
    \item[\textit{Characteristica universalis}] Universal formal language.
    \item[\textit{Calculus Ratiocinator}] Logical calculation framework. 
\end{description}
The objective was that: ``If two friends should come to a disagreement, it would suffice for them to take their pencils in their hands, sit down at the abacus, and say to each other - Let us calculate!''

\begin{definition}{Proof System}{proof-system}
    To prove a statement within a system, however complex, you have to:
    \begin{enumerate}
        \item Assert simple and obvious facts;
        \item Break the proof into a sequence of steps, each based on previous ones (chain of reasoning);
        \item Ensure each step is so simple and obvious that everyone would agree;
        \item Ensure the final step asserts the proposed statement.
    \end{enumerate}
    By doing so, everyone is ``forced'' to accept the statement.
\end{definition}

A proof system is composed of:
\begin{description}
    \item[Axioms] A list of basic facts so simple and self-evident that anyone agrees on them. 
    \item[Theorems] A dynamic list of known true facts.
    \item[Deduction Rules] Rules that combine known true facts to generate more true facts to be added to the theorem list.
        \subitem Example: \textit{Modus Ponens} --- if fact \(A\) is true and the fact \(A\Rightarrow B\) is true, then \(B\) is true.
\end{description}

Frege's axioms are very powerful: they can generate all of mathematics, and then something bad!
This is because \textit{Bertrand Russell} proved that \textbf{contradictions} arise from the system.

\begin{definition}{Unrestricted Comprehension Principle}{ucp}
    For every well-defined property \(P(\cdot)\), there is the set
    \[
        \{x\mid P(x)\}
    \]
    of all and only the objects having that property.
\end{definition}

As an example: if \(P(x)\equiv\) ``\(x\) is green'', then \(G=\{x\mid P(x)\}\) is the set of everything that is green. Now let \(P(x)\equiv\) ``\(x\) is a set''. Then \(S = \{x\mid P(x)\}\) is the set of all sets. Now \(G\) and \(S\) are sets and therefore elements of \(S\): \(G\in S, S\in S\). But on the other hand, neither \(G\) nor \(S\) are green, so \(G\notin G\) and \(S\notin G\).

\begin{theorem}{Russell's Paradox}{russell}
    Every set theory that contains an unrestricted comprehension principle leads to a contradiction.
\end{theorem}
\begin{proof}
    If \(x\) is a set, \(x\notin x\) (\(x\) does not belong to itself) is a well-defined property. By the unrestricted comprehension principle, the set of all sets that do not belong to themselves exists:
    \[
        R=\{x\mid x\notin x\}
    \]
    This introduces a contradiction, because \(R\in R \iff R\notin R\).
\end{proof}
This paradox is also known as the ``Barber Paradox''.

\noindent Ruling out unrestricted comprehension:
\begin{itemize}
    \item Russell and Whitehead, \gls{pm}
    \item Zermelo and Fraenkel set theory (\gls{zf})
\end{itemize}
A property \(P(x)\) can only be used to ``filter'' an already defined set. The notation:
\[
    \cancel{\{x\mid P(x)\}}
\] 
is now illegal; only \(\{x\in A\mid P(x)\}\) is writable, where \(A\) \textbf{MUST} be already defined and in the domain of \(P\).

\paragraph{Problem \# 2} Prove that the axioms of arithmetic are consistent. Both \gls{pm} and \gls{zf} prevent self-referential definitions, and therefore avoid paradoxical objects such as Russell's set \(R\), but more subtle inconsistencies might lurk inside them.

Indeed, the capability of formalizing integer arithmetic (possessed by both \gls{pm} and \gls{zf}) is enough to enable self-reference in a new way:

\paragraph{Kurt Gödel - \textit{Incompleteness theorems}} Reduction of logic to arithmetic: define a way to encode all of \gls{pm}'s logical statements as integer numbers and logical deduction rules as arithmetic calculations.

This showed a way to construct a number \(G\) corresponding to the statement: ``The statement represented by number \(G\) cannot be proved'':
\begin{itemize}
    \item If it is true, then by its own definition it has no proof, and so \gls{pm} is \textbf{incomplete} because there is at least one true statement that \gls{pm} cannot prove.
    \item If the statement is false, then the converse is true, meaning that a \textbf{false} statement can be proven, so by deduction it is true. This means that \gls{pm} is \textbf{inconsistent}. 
\end{itemize}
If we have to choose, we would prefer an incomplete system over an inconsistent system.

\begin{concept}{Hilbert's Entscheidungsproblem}{hilb-problem}
    Is there an automated, unambiguous procedure to decide whether a given statement is valid using the rules of logic?
\end{concept}

The answer of Alonzo Church and Alan Turing is:
\begin{center}
    \textbf{Nope!}
\end{center}

The following example, unlike Gödel's example, is not pathological; it has an actual impact in \gls{cs}.

\begin{algorithm}[H]
    \caption{Will it Halt?}
    \begin{algorithmic}
        \State \(n\leftarrow 1\) \Comment{Start from the first odd number}
        \Loop
            \State \(sum \leftarrow 0\)
            \For{\(i\leftarrow 1 \dots n-1\)} \Comment{Sum all divisors of \(n\)}
                \If{\(n\equiv 0 \pmod{i}\)} 
                    \State \(sum\leftarrow sum + i\)
                \EndIf
            \EndFor
            \If{\(sum = n\)} \Comment{If \(n\) is perfect, Halt!}
                \State \textbf{halt}
            \EndIf
            \State \(n\leftarrow n+2\) \Comment{Else move onto the next \textit{odd} number}
        \EndLoop
    \end{algorithmic}
\end{algorithm}

\begin{definition}{Perfect Number}{perfect-number}
    A natural number is \textit{perfect} if it is equal to the sum of its proper divisors.
\end{definition}

To create a program for which it is unknown whether it will halt or not, we can:
\begin{enumerate}
    \item Take an unproven conjecture in the form of ``All \(x\) are \(y\)''
    \item Run a program that systematically looks for a counterexample
    \item \dots
    \item Profit!
\end{enumerate}

In \gls{cs} we ``turn the tables'': maybe there exists a general procedure that, given an algorithm \(A\), tells us whether \(A\) will \textbf{halt}. Just running the program will not provide an answer (if it runs forever), but there should be some static analysis technique that looks at the algorithm and tells us if it will halt or not.

Suppose that there is a function \texttt{HALT(\(A\))} that takes \(A\), an algorithm, and \textbf{always} returns a correct answer:
\begin{align}
    \texttt{HALT}(A)=\begin{cases}
        \text{true} & \text{if } A \text{ halts}\\
        \text{false} & \text{if } A \text{ does not halt}
    \end{cases}
\end{align}

The contradiction:
\begin{algorithm}[H]
    \caption{Halting paradox}
    \begin{algorithmic}
        \Procedure{$R$}{}
        \If{\textsc{HALT}$(R)$}
            \State \textbf{loop}
            \Comment{If the HALT function predicts I am going to halt...}
        \Else \Comment{\textit{... then run forever.}}
            \State \textbf{halt}
            \Comment{If HALT predicts that I am running forever.}
        \EndIf \Comment{\textit{... then halt.}}
    \EndProcedure
    \end{algorithmic}
\end{algorithm}

This algorithm is self-referential and it isn't written in a well-defined programming language.

\chapter{Computability}
\section{Basic Definitions}

\begin{definition}{Alphabet}{alphabet}
    An alphabet is a set with a finite number of symbols.
\end{definition}

\noindent Usually, an alphabet (\cref{def:alphabet}) is denoted by \(\Sigma\). Now, let us compose words. To do so, we shall use the Cartesian product. As an example, the set of pairs of symbols will be \(\Sigma\times\Sigma=\Sigma^2\). The cardinality of this set is given by the product of the cardinalities of the composing sets: \(|\Sigma\times\Sigma|=|\Sigma|\cdot|\Sigma|\). To represent words with \(n\) symbols, we define the general set of length \(n\) by induction:
\[
    \Sigma^n=\begin{cases}
        \{\varepsilon\} & n=0\\
        \Sigma^{n-1}\times \Sigma & n>0
    \end{cases}
\]
This represents ``The set of all strings on \(\Sigma\) of length \(n\)''.

\noindent To define the set of all words, independent of their length, we take the union of all \(\Sigma^n\) sets:
\[
    \Sigma^0 \cup\Sigma^1 \cup\Sigma^2 \cup\ldots=\bigcup_{n\in\mathbb{N}}\Sigma^n=\Sigma^*
\]

\begin{definition}{Strings on an Alphabet}{strings}
    The strings on an alphabet \(\Sigma\) (\cref{def:alphabet}) are all the strings in its Kleene closure, which is the set:
    \[\Sigma^*=\bigcup_{n\in\mathbb{N}}\Sigma^n\]
\end{definition}

\begin{definition}{Enumeration}{enumeration}
    An enumeration of the strings (\cref{def:strings}) on an alphabet (\cref{def:alphabet}) is a 1:1 and onto function (bijection) that maps the set of natural numbers \(\mathbb{N}\) to the set \(\Sigma^*\).
    \[
        f:\mathbb{N}\rightarrow \Sigma^*
    \]
\end{definition}

\noindent By defining such a bijection between integers and strings, we establish that the set \(\mathbb{N}\) and the set \(\Sigma^*\) have the same cardinality (i.e., they are countably infinite).

\begin{definition}{Decision Function}{decision-function}
    A decision function is a function that takes a string (\cref{def:strings}) in the alphabet \(\Sigma^*\) and maps it to the set \(\{0,1\}\), representing TRUE and FALSE.
    \[
        f:\Sigma^*\rightarrow\{0,1\}
    \]
\end{definition}

\begin{definition}{Partial and Total Functions}{part-total-func}
    A \textbf{total function} maps every element of the domain to exactly one element of the codomain. A \textbf{partial function} is a function that is not total (i.e., it is undefined for at least one element in the domain).
\end{definition}

\begin{definition}{Equal and Distinct Functions}{equal-func}
    Two functions \(f\) and \(g\) are evaluated as equal or distinct based on the following criteria:
    \begin{align}
        f &= g \iff \forall s \in \Sigma^*, \; f(s) = g(s) \\
        f &\neq g \iff \exists s \in \Sigma^*, \; f(s) \neq g(s)
    \end{align}
\end{definition}

\begin{theorem}{Uncountability of Decision Functions}{uncount-dec}
    There is no enumeration (\cref{def:enumeration}) of decision functions (\cref{def:decision-function}).
    \begin{proof}
        Suppose, for the sake of contradiction, that there is a bijection mapping \(\mathbb{N}\) to the set of all decision functions (\cref{def:decision-function}):
        \[
            \mathbb{N}\rightarrow\{f:\Sigma^*\rightarrow \{0,1\}\}
        \]
        Because \(\Sigma^*\) is isomorphic to \(\mathbb{N}\) (as established by \cref{def:enumeration}), this implies a surjective mapping exists directly from any string \(s \in \Sigma^*\) to a decision function \(f_s\):
        \[
            s \mapsto f_s
        \]
        Since this mapping is assumed to be surjective, every possible decision function \(g\) must be accounted for in this enumeration:
        \[
            \forall g \in \{f:\Sigma^*\rightarrow\{0,1\}\}, \quad \exists s \in \Sigma^* \text{ such that } f_s = g
        \]
        Now, let us define a new decision function \(h\) (using Cantor's diagonal argument):
        \[
            h(s) = 1 - f_s(s) \quad \forall s \in \Sigma^*
        \]
        By this definition, it is mathematically true that for every string \(s\):
        \[
            h(s) \neq f_s(s) \implies h \neq f_s
        \]
        Thus, \(h\) is distinct from every \(f_s\) in our enumeration, meaning the function \(h\) is unmapped. This contradicts the assumption that our mapping is surjective. Therefore, a 1:1 mapping (enumeration) of decision functions cannot exist.
     \end{proof}
\end{theorem}
\noindent Thus, there is a fundamental difference in cardinality between the set of strings \(\Sigma^*\) (\cref{def:strings}) and the set of all decision functions \(\{f:\Sigma^*\rightarrow \{0,1\}\}\) (\cref{def:decision-function}). This distinction perfectly illustrates the boundary between \textbf{countable} and \textbf{uncountable} infinite sets.

\begin{definition}{Countable Set}{countable}
    A set \(S\) is \textbf{countable} (or countably infinite) if and only if there exists a bijection (a 1:1 and onto total function) from the set itself to the set of natural numbers \(\mathbb{N}\). Formally, there exists:
    \[
        \exists f: S \rightarrow \mathbb{N} \quad \text{such that } f \text{ is bijective (meaning } f^{-1}: \mathbb{N} \rightarrow S \text{ exists)}.
    \]
\end{definition}

\noindent Based on \cref{def:enumeration} and \cref{def:countable}, we can formally state that the set of strings \(\Sigma^*\) is countable, whereas the set of all decision functions is uncountable.
\begin{definition}{Language}{language}
    A language is a subset, defined by a decision function, of the set of all strings. \(L\subseteq\Sigma^*\) 
\end{definition}
\noindent Furthermore, there is a 1:1 correspondence (bijection) between the set of all decision functions and the power set (the set of all subsets) of \(\Sigma^*\), denoted as \(\mathcal{P}(\Sigma^*)\). 

\begin{definition}{Characteristic Function}{char-func}
    For any subset \(S \subseteq \Sigma^*\) (where \(S \in \mathcal{P}(\Sigma^*)\)), its \textbf{characteristic} (or indicator) \textbf{function} \(\chi_S\) is defined as:
    \[
    \begin{aligned}
        \chi_{S}: \Sigma^* &\rightarrow \{0,1\}\\
        w &\mapsto \begin{cases}
            1 & \text{if } w \in S\\
            0 & \text{if } w \notin S
        \end{cases}
    \end{aligned}
    \]
\end{definition}

\begin{theorem}{Bijection between Decision Functions and the Power Set}{bij-dec-power}
    There is a 1:1 correspondence (bijection) between the set of all decision functions on \(\Sigma^*\) and the power set \(\mathcal{P}(\Sigma^*)\).
    \begin{proof}
        Every decision function \(f\) defines a unique subset of \(\Sigma^*\). We can define a mapping \(C\) from the set of decision functions to the power set \(\mathcal{P}(\Sigma^*)\):
        \[
        \begin{aligned}
            C: \{f : \Sigma^* \rightarrow \{0,1\}\} &\rightarrow \mathcal{P}(\Sigma^*)\\
            f &\mapsto \underbrace{\{s \in \Sigma^* \mid f(s) = 1\}}_{f^{-1}(\{1\})}
        \end{aligned}
        \]
        Conversely, we can define an inverse mapping \(\chi\) using the characteristic function (\cref{def:char-func}) from the power set back to the set of decision functions:
        \[
        \begin{aligned}
            \chi: \mathcal{P}(\Sigma^*) &\rightarrow \{f : \Sigma^* \rightarrow \{0,1\}\}\\
            S &\mapsto \chi_S
        \end{aligned}
        \]
        Because \(C\) and \(\chi\) are inverses of each other (\(C(\chi_S) = S\) and \(\chi_{C(f)} = f\)), they form a bijection. 
    \end{proof}
\end{theorem}

\noindent \Cref{thm:uncount-dec} is a specific application of a much broader mathematical rule, known as Cantor's Theorem:

\begin{theorem}{Cantor's Theorem (General)}{set-pwr-set}
    For any set \(A\), the cardinality of its power set \(|\mathcal{P}(A)|\) is strictly greater than the cardinality of the set itself \(|A|\):
    \[
        |\mathcal{P}(A)| > |A|
    \]
    Furthermore, if \(A\) is a finite set with a cardinality of \(m\) (\(|A| = m\)), then the cardinality of its power set is exactly \(2^m\):
    \[
        |\mathcal{P}(A)| = 2^m
    \]
\end{theorem}

\noindent As an example of defining a set using a finite number of symbols, consider the set of prime numbers \(\mathbb{P}\):
\[
    \mathbb{P} = \{n \in \mathbb{N} \mid n > 1 \land \forall m \in \mathbb{N},\; 1 < m < n \implies m \nmid n\}
\]
An infinite set that can be unambiguously described using a finite string of symbols from a given alphabet is called \textbf{definable}.

\noindent Now, a question arises: \textbf{How many definable sets are there?} \\
They are infinite in number, but strictly \textbf{countable} (since their finite definitions are just strings, and the set of all strings is countable, as established in \cref{def:strings}).

\begin{concept}{Computable Language (Approximate Definition)}{comp-lang-approx}
    We can give an approximate definition of a ``computable language''. A language \(L \subseteq \Sigma^*\) is computable if there exists a program \(P\) (e.g., written in Python, C, etc.) such that for any string \(s\):
    \[
        P(s) = \begin{cases}
            1 & \text{if } s \in L \\
            0 & \text{if } s \notin L
        \end{cases}
    \]
    \textit{Note: The rigorous, formal definition (using Turing Machines) will be introduced later in the course.}
\end{concept}

\begin{definition}{Definable Language}{definable-language}
    A language \(L\) is \textbf{definable} if there exists a first-order predicate \(P\) such that:
    \[
        P(s) \iff s \in L
    \]
\end{definition}

\begin{observation}{Computability vs. Definability}{comp-vs-def}
    \textbf{All computable languages are definable.} The finite source code of the program that computes the language serves as its definition. 
    
    However, \textbf{not all definable languages are computable.} There are languages that can be perfectly defined using formal mathematical logic (such as the Halting Problem) for which no algorithm or program can ever exist.
\end{observation}
\section{A Computational Model: The Turing Machine}

\begin{observation}{Formulas as Strings}{form-strings}
    As a foundational premise: every mathematical formula, no matter how complex, can be unambiguously represented by a single finite string (\cref{def:strings}) of symbols.
\end{observation}

\noindent A program or algorithm can be represented and executed using a theoretical computational model known as a \textbf{Turing Machine}. The machine operates on a one-dimensional tape divided into discrete cells. 

Different definitions of the Turing Machine exist in literature:
\begin{itemize}
    \item \textbf{Tape bounds:} Some definitions require a tape with a defined starting point (e.g., a cell \(0\) exists, making it one-way infinite), while others define the tape as infinite in both directions.
    \item \textbf{Multiple tapes:} Some definitions allow for \(n\)-tapes. However, everything that can be computed on an \(n\)-tape machine can be simulated on a \(1\)-tape machine. Adding more tapes does not increase the computational power (or the class of computable functions) of the machine.
\end{itemize}

\noindent To process numbers and inputs effectively, we designate a \textbf{blank symbol} (usually denoted as \(\sqcup\)) indicating an empty cell, alongside the standard symbols from our alphabet \(\Sigma\) (\cref{def:alphabet}). 

A strict constraint of the machine is that it can only read one cell at a time and has a finite internal memory. To overcome the memory limitation, the machine leverages the infinite tape for storage during computation, which does not reduce its computational capability.

\begin{definition}{Turing Machine States and Transition Function}{tm-trans-func}
    The finite internal memory of the machine is represented by a set of \textbf{states}, denoted as \(Q\). There exists a specific initial state \(q_0 \in Q\) where the computation begins. 

    The computation is driven by a \textbf{transition function} \(F\), which maps the current state and the symbol currently being read to a new state, a symbol to write, and a directional head movement (\(\leftarrow\) for left, \(\rightarrow\) for right):
    \[
    \begin{aligned}
        F: \Sigma \times Q &\longrightarrow \Sigma \times Q \times \{\leftarrow, \rightarrow\}\\
        (\sigma, q) &\longmapsto (\sigma', q', \text{direction})
    \end{aligned}
    \]
    With this simple mechanism, we can represent every computable action, function, or calculation.
\end{definition}

\begin{observation}{The Halt State}{halt-state}
    There is a special state called \textsc{halt} that stops the execution of the transition function. Depending on the formalization, this is handled in the domain/codomain in two mathematically equivalent ways:
    \begin{itemize}
        \item If we consider \(\textsc{halt} \in Q\), the domain excludes it:
        \[
            F: \Sigma \times \left(Q \setminus \{\textsc{halt}\}\right) \longrightarrow \Sigma \times Q \times \{\leftarrow, \rightarrow\}
        \]
        \item If we consider \(\textsc{halt} \notin Q\), it is added to the codomain:
        \[
            F: \Sigma \times Q \longrightarrow \Sigma \times \left(Q \cup \{\textsc{halt}\}\right) \times \{\leftarrow, \rightarrow\}
        \]
    \end{itemize}
\end{observation}

\subsection{Example: Binary Incrementor}
Assume our alphabet (\cref{def:alphabet}) is \(\Sigma = \{\sqcup, 0, 1\}\), where \(\sqcup\) is the default blank symbol. We want to define a transition function (\cref{def:tm-trans-func}) that adds \(1\) to a binary number already written on the tape. 

We define our states as \(Q = \{R, A\}\) (abbreviating \(q_{\text{right}}\) and \(q_{\text{add\_one}}\)). The machine will first move to the far right of the string, and then move left while performing the binary addition (flipping \(1\)s to \(0\)s and carrying the \(1\) until it finds a \(0\) or a blank).

The transition function \(F\) (\cref{def:tm-trans-func}) can be represented using a state table. Note that when transitioning into the \textsc{halt} state (\cref{obs:halt-state}), the computation terminates immediately:

\begin{center}
    \renewcommand{\arraystretch}{1.5}
    \begin{tabular}{r | c c c}
        \textbf{State \textbackslash\ Read} & \(\sqcup\)& \(\mathbf{0}\) & \(\mathbf{1}\) \\
        \hline
        \(R\) \textit{(move\_right)} & \((\sqcup, A, \leftarrow)\) & \((0, R, \rightarrow)\) & \((1, R, \rightarrow)\) \\
        \(A\) \textit{(add\_one)} & \((1, \textsc{halt}, \leftarrow)\) & \((1, \textsc{halt}, \leftarrow)\) & \((0, A, \leftarrow)\) \\
    \end{tabular}
\end{center}

\noindent Nowadays, automata are frequently represented by a directed graph where nodes represent the states (\cref{def:tm-trans-func}) and directed, labeled arcs represent the transitions. The standard notation for an arc label in a Turing Machine is \(\text{Read} / \text{Write}, \text{Direction}\). 

Thus, our binary incrementor can be visually represented by the following state diagram:

\begin{figure}[H]
    \centering
    \begin{tikzpicture}[
        >=Stealth, 
        shorten >=1pt, 
        auto, 
        node distance=4.5cm, 
        semithick,
        every state/.style={thick, fill=gray!10}
    ]
        % Nodes
        \node[state, initial, initial text=start] (R) {\(R\)};
        \node[state] (A) [right=of R] {\(A\)};
        \node[state, accepting] (H) [right=of A] {\textsc{halt}};

        % Edges
        \path[->]
        (R) edge [loop above] node[align=center] {\(0 / 0, \rightarrow\) \\ \(1 / 1, \rightarrow\)} (R)
            edge              node {\(\sqcup / \sqcup, \leftarrow\)} (A)
        (A) edge [loop above] node {\(1 / 0, \leftarrow\)} (A)
            edge              node[align=center] {\(0 / 1, \leftarrow\) \\ \(\sqcup / 1, \leftarrow\)} (H);
    \end{tikzpicture}
    \caption{Incrementor state graph}
    \label{fig:incrementor-graph}
\end{figure}

\backmatter
\chapter{Appendix}

% Generate automated lists of definitions, concepts, theorems, and axioms
\tcblistof[\section]{deflist}{List of Definitions}
\tcblistof[\section]{cptlist}{List of Concepts}
\tcblistof[\section]{thmlist}{List of Theorems}
\tcblistof[\section]{axlist}{List of Axioms}

\clearpage
% Set a nice gap between the two columns
\setlength{\columnsep}{1cm}

% Span the section title across both columns
\begin{multicols}{2}[\section{List of Acronyms}]
    \raggedright
    % Disable the default glossary title so it doesn't print twice
    \renewcommand{\glossarysection}[2][]{}
    \printglossary[type=\acronymtype]
\end{multicols}

\end{document}